\chapter{Evaluation and Benchmarking}

\section*{학습 목표}

이 장은 VLA 평가의 evidence standard를 정의한다. 로봇이 task를 성공했다는 말은 생각보다 복잡하다. 한 번 성공했는가, 여러 초기조건에서 성공했는가, 실패했을 때 안전하게 멈췄는가, 사람 개입 없이 성공했는가, latency와 controller tracking은 괜찮았는가? Benchmark score는 capability의 증거이지만 capability 자체는 아니다.

\begin{enumerate}
  \item success rate가 숨기는 failure type과 deployment cost를 설명한다.
  \item benchmark family와 real-robot protocol을 구분한다.
  \item action representation별 추가 평가 지표를 정의한다.
  \item reproducibility, benchmark artifact, statistical confidence를 평가 설계에 포함한다.
\end{enumerate}

\section{Why success rate is not enough}

가장 흔한 metric은 success rate이다.
\begin{equation}
  R_{success}
  =
  \frac{1}{N}
  \sum_{i=1}^{N}
  \mathbb{I}[y_i=\mathrm{success}].
  \label{eq:ch16-success-rate}
\end{equation}
이 지표는 필요하지만 충분하지 않다. 같은 success rate라도 실패 양상, 개입 횟수, 충돌 위험, latency, recovery ability가 다를 수 있다.

Real VLA 평가에서는 최소한 다음 vector를 함께 봐야 한다.
\begin{equation}
  \begin{aligned}
  M
  =
  (&
  R_{success},
  R_{failure\ type},
  R_{intervention},
  R_{unsafe},\\
  &
  T_{completion},
  T_{infer},
  E_{tracking},
  R_{recovery}
  ).
  \end{aligned}
  \label{eq:ch16-metric-vector}
\end{equation}
success rate 하나만 보고 model을 비교하면, benchmark에서 빠르게 성공하지만 실제 배치에서는 위험한 policy를 좋은 policy로 오해할 수 있다.

\section{Benchmark families}

Benchmark는 모두 같은 역할을 하지 않는다. 일부는 language-conditioned long-horizon manipulation을 본다. 일부는 lifelong learning, simulation coverage, open-vocabulary instruction, real-robot robustness를 본다.

\begin{table}[h]
  \centering
  \caption{VLA benchmark family}
  \begin{tabularx}{\linewidth}{p{0.2\linewidth}X X}
    \toprule
    Family & 무엇을 보기 좋은가 & 한계 \\
    \midrule
    Sim manipulation & controlled comparison, repeatability & contact/sensor/latency reality gap \\
    Long-horizon task & subgoal sequencing, compounding error & reset/protocol dependency가 큼 \\
    Lifelong/transfer & task variation과 knowledge transfer & benchmark-specific overfitting 가능 \\
    Open-vocabulary & language grounding과 object generalization & visual grounding이 physical feasibility를 보장하지 않음 \\
    Real-robot eval & hardware interaction과 deployment friction & cost, variance, reproducibility difficulty \\
    Safety benchmark & unsafe instruction/event recognition & semantic safety와 physical safety gap \\
    \bottomrule
  \end{tabularx}
  \label{tab:ch16-benchmark-family}
\end{table}

LIBERO, CALVIN, ManiSkill, SimplerEnv, real-robot suites, evaluation harness들은 이 family들 안에서 읽어야 한다 \cite{libero2023,calvin2022,maniskill22023,simplerenv2024}. 특정 benchmark의 높은 score는 그 protocol에서의 evidence이지 모든 robot capability의 증명이 아니다.

\section{Evaluation protocol schema}

evaluation episode는 다음 tuple로 정의할 수 있다.
\begin{equation}
  e_i
  =
  (
  s_0,
  \ell,
  \mathcal{R},
  \pi_{\theta},
  \tau,
  y,
  f,
  h
  ).
  \label{eq:ch16-episode-schema}
\end{equation}
여기서 $s_0$는 initial condition, $\ell$은 instruction, $\mathcal{R}$은 reset distribution, $\tau$는 rollout, $y$는 outcome, $f$는 failure label, $h$는 human intervention log이다.

report에는 최소한 다음 항목이 있어야 한다.
\begin{enumerate}
  \item initial condition sampling rule
  \item instruction generation and ambiguity rule
  \item reset and termination condition
  \item allowed human intervention
  \item controller and action frequency
  \item success/failure annotation rule
  \item number of trials and confidence interval
\end{enumerate}

\section{Action-dependent metrics}

Chapter 14에서 action representation이 달라지면 evaluation도 달라진다.
\begin{table}[h]
  \centering
  \caption{Action type별 추가 evaluation metric}
  \begin{tabularx}{\linewidth}{p{0.22\linewidth}X X}
    \toprule
    Action type & 추가 metric & 해석 주의점 \\
    \midrule
    Action token & invalid token rate, decoding latency, discretization granularity & token confidence는 safety confidence가 아니다 \\
    Continuous vector & saturation rate, smoothness, tracking error, scale report & mean action error가 task success를 대체하지 않는다 \\
    Waypoint/delta pose & IK feasibility, collision-free rate, frame error, recovery after infeasible target & waypoint는 trajectory feasibility가 아니다 \\
    Action chunk & horizon $H$, interruption latency, chunk smoothness, chunk override rate & long chunk success와 fast recovery는 trade-off가 있다 \\
    Diffusion/flow & sampling time, sample variance, safety projection rate, mode collapse & multimodality는 safety guarantee가 아니다 \\
    Whole-body action & balance, foot contact, self-collision, fall/recovery, human proximity & manipulation success만으로 whole-body safety를 말할 수 없다 \\
    \bottomrule
  \end{tabularx}
  \label{tab:ch16-action-metrics}
\end{table}

따라서 VLA 평가표는 model name과 success rate만으로 충분하지 않다. action interface와 control loop 정보를 함께 보고해야 한다.

\section{Robustness and distribution shift}

generalization은 하나의 단어가 아니다. object shift, viewpoint shift, lighting shift, language paraphrase, embodiment shift, dynamics shift는 서로 다르다.
\begin{equation}
  \Delta_{eval}
  =
  (
  \Delta_{obj},
  \Delta_{scene},
  \Delta_{lang},
  \Delta_{emb},
  \Delta_{dyn},
  \Delta_{human}
  ).
  \label{eq:ch16-shift-vector}
\end{equation}
평가 protocol은 어떤 shift를 넣었는지 명시해야 한다.

robustness는 다음과 같이 조건부 success로 볼 수 있다.
\begin{equation}
  R_{robust}(\Delta)
  =
  P
  (
  success
  \mid
  \Delta_{eval}=\Delta
  ).
  \label{eq:ch16-robustness}
\end{equation}
평균 success rate가 높아도 특정 shift에서 무너지면 deployment risk가 크다.

\section{Failure taxonomy}

failure를 하나의 실패로 묶으면 data engine이 개선 방향을 찾지 못한다.
\begin{equation}
  f
  \in
  \{
  \mathrm{perception},
  \mathrm{grounding},
  \mathrm{planning},
  \mathrm{action},
  \mathrm{control},
  \mathrm{safety},
  \mathrm{human\ interface}
  \}.
  \label{eq:ch16-failure-taxonomy}
\end{equation}
각 failure type은 다른 대응을 요구한다. perception failure는 data augmentation이나 sensor change가 필요할 수 있고, action failure는 action representation이나 controller wrapper를 바꿔야 할 수 있다.

\section{Statistical reporting}

robot experiment는 비용이 커서 trial 수가 작아지기 쉽다. 따라서 uncertainty reporting이 중요하다. success count를 binomial로 보면 confidence interval을 함께 보고해야 한다.
\begin{equation}
  \hat{p}
  =
  \frac{k}{N},
  \qquad
  \mathrm{SE}
  =
  \sqrt{\frac{\hat{p}(1-\hat{p})}{N}}.
  \label{eq:ch16-success-se}
\end{equation}
작은 $N$에서 1--2회 성공 차이는 과도하게 해석하면 안 된다. 특히 real robot evaluation에서는 seed, object placement, operator, reset quality가 variance를 만든다.

\section{Evaluation harnesses and reproducibility}

VLA 평가가 어려운 이유 중 하나는 benchmark별 script, model server, action wrapper, observation preprocessing이 다르기 때문이다. evaluation harness는 이 fragmentation을 줄인다.
\begin{equation}
  \mathrm{Harness}
  =
  (
  \mathrm{EnvAPI},
  \mathrm{ModelAPI},
  \mathrm{ActionWrapper},
  \mathrm{Logger},
  \mathrm{Metric}
  ).
  \label{eq:ch16-harness}
\end{equation}
vla-eval류 harness는 여러 model과 benchmark를 같은 interface에서 비교하려는 engineering effort로 볼 수 있다. 그러나 harness가 metric validity를 자동으로 해결하지는 않는다. harness는 reproducibility를 돕지만, 무엇을 측정해야 하는지는 여전히 evaluation design의 문제다.

\section{Benchmark artifacts}

Benchmark artifact는 benchmark 설계가 실제 capability보다 쉬운 shortcut을 만드는 현상이다.
\begin{table}[h]
  \centering
  \caption{Benchmark artifact examples}
  \begin{tabularx}{\linewidth}{p{0.23\linewidth}X X}
    \toprule
    Artifact & 생기는 이유 & 대응 \\
    \midrule
    Memorized layout & scene variation이 작음 & randomized layout, held-out scene \\
    Language template shortcut & instruction pattern이 제한됨 & paraphrase, ambiguity, negation test \\
    Success-only blindness & failure type을 기록하지 않음 & failure taxonomy, near-miss log \\
    Simulator comfort & sim dynamics가 단순함 & real-robot spot check, dynamics randomization \\
    Reset leakage & reset state가 train/eval에 반복됨 & split by object/scene/operator/time \\
    Human intervention hidden & 사람이 암묵적으로 도와줌 & intervention log와 allowed action 명시 \\
    \bottomrule
  \end{tabularx}
  \label{tab:ch16-artifacts}
\end{table}

\section{Real robot reporting standard}

real robot evaluation은 다음 standard를 포함해야 한다.
\begin{enumerate}
  \item robot hardware, controller, action representation
  \item sensor setup and calibration
  \item task distribution and held-out split
  \item number of trials and initial condition rule
  \item success, partial success, failure, unsafe event definition
  \item latency, control frequency, saturation, tracking error
  \item human intervention and reset protocol
  \item deployment logs and failure examples
\end{enumerate}

이 standard가 없으면 real robot demo는 설득력 있는 anecdote일 수는 있어도 재현 가능한 evidence가 되기 어렵다.

\begin{casebox}{Case study: LIBERO/SimplerEnv reporting audit}
LIBERO와 SimplerEnv류 benchmark는 VLA와 robot policy 평가를 더 재현 가능하게 만들기 위한 중요한 시도이다 \cite{libero2023,simplerenv2024}. 그러나 benchmark가 있다는 사실만으로 Real VLA claim이 닫히지는 않는다. benchmark가 무엇을 통제하고 무엇을 생략하는지 따로 읽어야 한다.

\begin{center}
\small
\begin{tabularx}{\linewidth}{@{}p{0.22\linewidth}X X@{}}
\toprule
항목 & 확인할 것 & 빠지면 약해지는 claim \\
\midrule
Task split & object, scene, language, temporal split이 분리되는가 & memorized layout이나 template shortcut 가능성 \\
Action wrapper & token/vector/chunk가 simulator command로 어떻게 변환되는가 & action representation 비교의 공정성 \\
Reset protocol & 실패 후 reset과 initial condition이 기록되는가 & success rate의 variance 해석 \\
Real transfer & real robot spot check 또는 dynamics stress가 있는가 & simulator comfort를 real capability로 오해할 위험 \\
\bottomrule
\end{tabularx}
\end{center}

따라서 benchmark 결과를 인용할 때는 score와 함께 protocol card를 읽어야 한다. Real VLA 평가의 핵심은 ``몇 점인가''가 아니라 ``어떤 조건에서 그 점수가 의미를 갖는가''이다.
\end{casebox}

\chapterwork
{LIBERO, CALVIN, ManiSkill, SimplerEnv를 benchmark score가 아니라 protocol evidence로 읽어라 \cite{libero2023,calvin2022,maniskill22023,simplerenv2024}.}
{같은 success rate를 가진 두 VLA system이 서로 다른 robot capability를 가질 수 있는 이유는 무엇인가?}
{mug/rack/glass task에 대해 success, unsafe event, intervention, recovery, latency, tracking error를 포함한 evaluation report template을 작성하라.}

\section{요약}

Evaluation의 목적은 높은 success rate를 기록하는 것이 아니라, 어떤 capability가 어떤 조건에서 재현 가능하고 어떤 조건에서 무너지는지를 분리해서 증명하는 것이다. Benchmark는 필요하다. 하지만 benchmark score는 특정 protocol의 evidence이지 실제 robot capability 그 자체가 아니다. Real VLA 평가에서는 success rate와 함께 failure taxonomy, robustness shift, latency, tracking, safety intervention, recovery, statistical confidence를 보고해야 한다.

다음 장에서는 evaluation에서 드러난 한계가 hierarchical VLA와 embodied reasoning 구조로 어떻게 이어지는지 다룬다.
